\clearpage\SourceLeaf{116}
\pdfbookmark[1]{Göttingische gelehrte Anzeigen. 1816 März 2.}{anzeige1816}
\begin{center}
\rule{0.45\linewidth}{0.4pt}\par
\textbf{Göttingische gelehrte Anzeigen.\quad 1816 März 2.}\par
\rule{0.45\linewidth}{0.4pt}
\end{center}

Nachdem die in diesen Anzeigen vom v. J. angezeigte Vorlesung des
Herrn Prof. \textsc{Gauss} bereits abgedruckt war, hatte derselbe bei
fortgesetzter Beschäftigung mit demselben Gegenstande das Glück,
dasselbe Ziel noch auf einem ganz neuen Wege zu erreichen. Er hat
hierüber am 30. Januar der königl. Societät eine kleine Abhandlung
eingereicht, die die Aufschrift führt:

\begin{center}
\emph{Theorematis de resolubilitate functionum algebraicarum integrarum
in factores reales demonstratio tertia.}
\end{center}

Die beiden frühern Beweise des wichtigsten Lehrsatzes in der Lehre
von den Gleichungen unterschieden sich dadurch, dass der erstere sehr
kurz und einfach, aber zum Theil auf geometrische Betrachtungen
gegründet war, der andere hingegen rein analytisch aber viel
complicirter, so dass es unmöglich wurde, in dem engen Raume dieser
Blätter einen genügenden Auszug daraus zu geben. Dagegen ist nun der
gegenwärtige dritte auf gänzlich verschiedenen Principien beruhende
Beweis ebenfalls rein analytisch, übertrifft aber selbst den ersten
so sehr an Einfachheit und Kürze, dass wir hier ganz füglich alles
Wesentliche desselben auf wenigen Seiten mittheilen können.

Es bezeichne $X$ die algebraische Function der unbestimmten Grösse $x$,
\[
x^m+Ax^{m-1}+Bx^{m-2}+Cx^{m-3}+\mathrm{u.\,s.\,w.}
\]
\SourceLeaf{117}
wo die Coëfficienten $A,B,C$ u.\,s.\,w. bestimmte reelle
Grössen sind. Es seien ferner $r$ und $\varphi$ zwei neue unbestimmte
Grössen, und man setze
\[
\begin{aligned}
&r^m\cos m\varphi+Ar^{m-1}\cos(m-1)\varphi
\\
&\quad+Br^{m-2}\cos(m-2)\varphi+Cr^{m-3}\cos(m-3)\varphi
+\mathrm{u.\,s.\,w.}=t\\
&r^m\sin m\varphi+Ar^{m-1}\sin(m-1)\varphi
\\
&\quad+Br^{m-2}\sin(m-2)\varphi+Cr^{m-3}\sin(m-3)\varphi
+\mathrm{u.\,s.\,w.}=u\\
&mr^m\cos m\varphi+(m-1)Ar^{m-1}\cos(m-1)\varphi
\\
&\quad+(m-2)Br^{m-2}\cos(m-2)\varphi+(m-3)Cr^{m-3}\cos(m-3)\varphi
+\mathrm{u.\,s.\,w.}=t'\\
&mr^m\sin m\varphi+(m-1)Ar^{m-1}\sin(m-1)\varphi
\\
&\quad+(m-2)Br^{m-2}\sin(m-2)\varphi+(m-3)Cr^{m-3}\sin(m-3)\varphi
+\mathrm{u.\,s.\,w.}=u'\\
&mmr^m\cos m\varphi+(m-1)^2Ar^{m-1}\cos(m-1)\varphi
\\
&\quad+(m-2)^2Br^{m-2}\cos(m-2)\varphi\\
&\quad+(m-3)^2Cr^{m-3}\cos(m-3)\varphi
+\mathrm{u.\,s.\,w.}=t''\\
&mmr^m\sin m\varphi+(m-1)^2Ar^{m-1}\sin(m-1)\varphi
\\
&\quad+(m-2)^2Br^{m-2}\sin(m-2)\varphi\\
&\quad+(m-3)^2Cr^{m-3}\sin(m-3)\varphi
+\mathrm{u.\,s.\,w.}=u''\\
&\frac{(tt+uu)(tt''+uu'')+(tu'-ut')^2-(tt'+uu')^2}
{r(tt+uu)^2}=y
\end{aligned}
\]



Der Factor $r$ kann offenbar aus dem Nenner des Ausdrucks für $y$
wegeschafft werden, da $t',u',t'',u''$ durch $r$ theilbar sind.

Es sei ferner $R$ eine beliebige bestimmte positive Grösse, mit der
einzigen Einschränkung, dass sie grösser sei als die grösste von folgenden
\[
mA\surd2,\quad\surd(mB\surd2),\quad
{}^3\!\surd(mC\surd2),\quad{}^4\!\surd(mD\surd2)\ \mathrm{etc.}
\]
abgesehen von dem Zeichen von $A,B,C$ u.\,s.\,w., oder indem man
alle als positiv betrachtet.

Endlich mögen noch folgende Bezeichnungen Statt finden:
\[
\begin{aligned}
&R^m\cos45^0+AR^{m-1}\cos(45^0+\varphi)
\\
&\quad+BR^{m-2}\cos(45^0+2\varphi)+CR^{m-3}\cos(45^0+3\varphi)
+\mathrm{u.\,s.\,w.}=T\\
&R^m\sin45^0+AR^{m-1}\sin(45^0+\varphi)
\\
&\quad+BR^{m-2}\sin(45^0+2\varphi)+CR^{m-3}\sin(45^0+3\varphi)
+\mathrm{u.\,s.\,w.}=U\\
&mR^m\cos45^0+(m-1)AR^{m-1}\cos(45^0+\varphi)
\\
&\quad+(m-2)BR^{m-2}\cos(45^0+2\varphi)
\\
&\quad+(m-3)CR^{m-3}\cos(45^0+3\varphi)+\mathrm{u.\,s.\,w.}=T'\\
&mR^m\sin45^0+(m-1)AR^{m-1}\sin(45^0+\varphi)
\\
&\quad+(m-2)BR^{m-2}\sin(45^0+2\varphi)
\\
&\quad+(m-3)CR^{m-3}\sin(45^0+3\varphi)+\mathrm{u.\,s.\,w.}=U'
\end{aligned}
\]

Nach allen diesen Vorbereitungen, welche wir der bequemern Uebersicht
wegen zusammengestellt haben, ergeben sich leicht nachstehende Folgerungen:

I. Die Grösse $T$ ist nothwendig positiv, welchen Werth man auch immer
der Grösse $\varphi$ beilege. Dies ist leicht zu übersehen, wenn man $T$
in folgende Form bringt:
\[
\begin{aligned}
&\frac{R^{m-1}}{m\surd2}(R+mA\surd2.\cos(45^{\mathrm{a}}+\varphi))\\
&\quad+\frac{R^{m-2}}{m\surd2}(RR+mB\surd2.\cos(45^0+2\varphi))\\
&\quad+\frac{R^{m-3}}{m\surd2}(R^3+mC\surd2.\cos(45^0+3\varphi))\\
&\quad+\frac{R^{m-4}}{m\surd2}(R^4+mD\surd2.\cos(45^0+4\varphi))\\
&\quad+\mathrm{u.\,s.\,w.}
\end{aligned}
\]
\SourceLeaf{118}
da jeder Theil einzeln genommen nothwendig positiv wird.
Auf ähnliche Weise ist leicht zu beweisen, dass für jeden Werth von
$\varphi$ auch $U,T',U'$ nothwendig positiv werden.

II. Für $r=R$ erhalten die Functionen $t,u,t',u'$ der Reihe nach
folgende Werthe
\[
\begin{gathered}
T\cos(45^0+m\varphi)+U\sin(45^0+m\varphi)\\
T\sin(45^0+m\varphi)-U\cos(45^0+m\varphi)\\
T'\cos(45^0+m\varphi)+U'\sin(45^0+m\varphi)\\
T'\sin(45^0+m\varphi)-U'\cos(45^0+m\varphi)
\end{gathered}
\]
Hieraus folgt, dass für $r=R$, $tt+uu=TT+UU$, $tt'+uu'=TT'+UU'$
also beide nothwendig positiv werden.

III. Es kann also für keinen Werth von $r$, der grösser ist als jede
der Grössen $mA\surd2$, $\surd(mB\surd2)$, ${}^3\!\surd(mC\surd2)$
u.\,s.\,w., zugleich $t=0$ und $u=0$ werden.

Das Haupttheorem ist nun folgendes:

\emph{Innerhalb der Grenzen $r=0$, $r=R$, und $\varphi=0$,
$\varphi=360^0$ (einschl.) gibt es gewiss Werthe für $r$ und $\varphi$,
aus denen zugleich $t=0$ und $u=0$ wird.}

Der Beweis davon wird auf folgende Art geführt. Wenn man annimmt,
das Theorem sei nicht wahr, so folgt, dass $tt+uu$ für jede Werthe
von $r$ und $\varphi$ zwischen den angegebenen Grenzen immer positiv,
und folglich $y$ immer \emph{endlich} werden muss. Betrachten wir nur
den Werth des doppelten Integrals
\[
\iint y\,\mathrm{d}r\,\mathrm{d}\varphi
\]
von $r=0$ bis $r=R$, und von $\varphi=0$ bis $\varphi=360^0$
erstreckt, der also eine ganz bestimmte endliche Grösse haben wird.
Dieser mit $\Omega$ zu bezeichnende Werth kann auf zweifache Art
gefunden werden, indem man entweder zuerst nach $\varphi$ und dann
nach $r$ integrirt, oder in umgekehrter Ordnung. Beide Wege müssen
nothwendig zu einerlei Resultate führen.

Nun hat man aber, wenn man $r$ als beständig betrachtet, unbestimmt
\[
\int y\,\mathrm{d}\varphi=\frac{tu'-ut'}{r(tt+uu)}
\]
wie man sich leicht durch Differentiation versichern kann. Eine beständige Grösse
\SourceLeaf{119}
ist bei der Integration nicht hinzuzufügen, insofern diese von
$\varphi=0$ anfangen soll, da für diesen Werth offenbar
\[
\frac{tu'-ut'}{tt+uu}=0
\]
wird: und da eben dies für $\varphi=360^0$ gilt, so ist der Werth von
$\int y\,d\varphi$, von $\varphi=0$ bis $\varphi=360^0$ erstreckt, $=0$,
bei jedem Werthe von $r$. Hieraus schliessen wir also $\Omega=0$.

Von der andern Seite hat man, wenn man $\varphi$ als beständig betrachtet,
unbestimmt
\[
\int y\,dr=\frac{tt'+uu'}{tt+uu}
\]
wie ebenfalls leicht durch die Differentiation nach $r$ bestätigt wird.
Auch hier ist keine beständige Grösse hinzuzusetzen, insofern die Integration
von $r=0$ anfangen soll. Das Integral $\int y\,dr$, von $r=0$ bis $r=R$
ausgedehnt, wird folglich nach II
\[
=\frac{TT'+UU'}{TT+UU}
\]
also gewiss positiv für jeden Werth von $\varphi$. Hieraus wird demnach
auch $\Omega$, d.\,i. der Werth des Integrals
\[
\int\frac{TT'+UU'}{TT+UU}\,d\varphi
\]
von $\varphi=0$ bis $\varphi=360^0$ ausgedehnt, nothwendig positiv werden,
welches mit dem erstern Resultate in Widerspruch steht. Die Voraussetzung
war folglich unstatthaft, und dadurch ist die Wahrheit des Theorems selbst
bewiesen.

Dass nun für jede Werthe von $r$ und $\varphi$, welche zugleich $t=0$
und $u=0$ geben,
\[
r(\cos\varphi+\sin\varphi\surd(-1))\quad\hbox{und}\quad
r(\cos\varphi-\sin\varphi\surd(-1))
\]
Wurzeln der Gleichung $X=0$ sind, oder dass die Function $X$ dann
entweder den Factor der zweiten Ordnung
\[
xx-2r\cos\varphi.x+rr
\]
(wenn weder $r$ noch $\sin\varphi$ verschwinden) oder den Factor der
ersten Ordnung
\[
x\mp r
\]
\SourceLeaf{120}
(wenn entweder $\sin\varphi=0$, also $\cos\varphi=\pm1$,
oder $r=0$ wird) enthält, ist bekannt genug, und braucht hier nicht erst
weiter entwickelt zu werden. Was von $X$ so eben bewiesen ist, gilt
nachher wieder von dem Quotienten, wenn $X$ durch jenen Factor dividirt
ist, u.\,s.\,f., woraus die vollständige Zerlegbarkeit der Function $X$ in
dergleichen Factoren erhellet.

Durch diese kurze Reihe von Schlüssen ist das eigentliche Ziel vollständig
erreicht. Nicht gefordert, aber doch gewünscht werden konnte noch eine
Erläuterung, in wiefern der obige Widerspruch --- der unvermeidlich war,
wenn man das Theorem als unwahr ansah --- wegfalle, wenn man von der
Wahrheit des Theorems ausgeht. Der Verf. hat darüber einige Winke gegeben,
die hier nur kurz berührt werden können. Ist das Theorem wahr, so sind
nicht mehr alle Werthe von $y$ endlich; es ist folglich nicht mehr ohne
Weiteres verstattet, $\iint y\,dr\,d\varphi$ als eine wirkliche bestimmte
Grösse anzunehmen, und man darf sich daher nicht wundern, dass der blosse
blinde Mechanismus des Calcüls, auf verschiedenen Wegen widersprechende
Resultate liefert. Die Analyse pflegt sehr häufig sich so zu helfen, dass
sie auf Fragen, die man ihr ohne Einschränkung vorlegt, obgleich sie in
gewissen Fällen ungereimt sein können, nur halbbestimmte Antworten gibt.
So ist es bei der Bestimmung des Werthes des Integrals $\iint y\,dr\,d\varphi$.
Soll dasselbe allgemein, von $r=k$ bis $r=l$ und von $\varphi=K$ bis
$\varphi=L$ erstreckt werden, und bezeichnet man den Werth von $\frac{u}{t}$
\[
\begin{array}{rcl}
\hbox{für }r=k,\quad\varphi=K&\hbox{durch}&\theta\\
r=k,\quad\varphi=L&\hbox{durch}&\theta'\\
r=l,\quad\varphi=K&\hbox{durch}&\theta''\\
r=l,\quad\varphi=L&\hbox{durch}&\theta'''
\end{array}
\]
so geben die analytischen Operationen für jenes Integral den Ausdruck
\[
\mathrm{Arc\,tg}\,\theta-\mathrm{Arc\,tg}\,\theta'
-\mathrm{Arc\,tg}\,\theta''+\mathrm{Arc\,tg}\,\theta'''
\]

Das Integral hat in der That nur dann einen wahren Werth, wenn $y$
innerhalb der angegebenen Grenzen immer endlich bleibt: dieser wahre
Werth ist dann unter obigem Ausdruck allerdings begriffen, aber an sich
dadurch noch nicht ganz bestimmt, da die Function Arc. tg. eine
vielförmige ist, und erst aus anderweitigen (übrigens nicht schwierigen)
Betrachtungen muss entschieden werden, \emph{welche}
\SourceLeaf{121}
Werthe dieser Function im bestimmten Fall zu nehmen sind.
So oft hingegen innerhalb der angegebenen Grenzen $y$ irgendwo unendlich
wird, ist eigentlich die Frage nach dem Werthe von $\iint y\,dr\,d\varphi$
ungereimt; unterfängt man sich, diese Ungereimtheit ignorirend, dennoch,
sie zu beantworten, so darf man sich nicht befremden lassen, auf einem
Wege diese, auf andern jene Antwort zu erhalten, welche verschiedene
Beantwortungen inzwischen allemal unter dem obigen halbbestimmten
Ausdrucke begriffen sind.

\par\medskip\noindent\rule{0.25\linewidth}{0.4pt}\par\smallskip
\begingroup\small
\emph{Handschriftliche Bemerkungen.}

Die \emph{Bedingung} für $R$ kann auch so gestellt werden, dass es
grösser ist als $1$ und als die absolute Summe der Grössen $A,B,C$ etc.
in $\surd2$ multiplicirt. Es sei diese absolute Summe
$=\alpha A+\beta B+\gamma C+\delta D+\mathrm{etc.}$ so dass jeder der
Coëfficienten $\alpha,\beta,\gamma$ etc. entweder $+1$ oder $-1$ wird.
Dann ist der positive Werth von $T$ u.\,s.\,w. aus folgender Form klar
\[
\begin{aligned}&T=R^{m-1}\cos45^0\bigl(R-\surd2.(\alpha A+\beta B+\gamma C+\mathrm{etc.})\bigr)\\
&\quad+R^{m-1}\bigl(\alpha A+A\cos(45^0+\varphi)\bigr)\\
&\quad+R^{m-2}\bigl(\beta BR+B\cos(45^0+2\varphi)\bigr)+\mathrm{etc.}
\end{aligned}
\]

\emph{Lehrsatz.} Sind $a,b,c\,.\,.\,.\,m,n$ die Wurzeln der Gleichung
$fx=0$, $a',b',c'\,.\,.\,.\,m'$ die Wurzeln der Gleichung $f'x=0$.
wo $f'x=\frac{dfx}{dx}$, und werden durch dieselben Buchstaben die
entsprechenden Punkte in plano bezeichnet, so ist, wenn man sich in
$a,b,c\,.\,.\,.\,m,n$ gleiche abstossende oder anziehende Massen denkt,
die im umgekehrten Verhältniss der Entfernung wirken, in
$a',b',c'\,.\,.\,.\,m'$ Gleichgewicht.
\endgroup

\par\medskip\begin{center}\rule{0.25\linewidth}{0.4pt}\end{center}
